\chapter{Fault-Tolerance Framework: DMR and TMR}
\label{ch:framework}

Building on the modular accelerator of Chapter~\ref{ch:architecture}, this
chapter presents the fault-tolerance framework. Section~\ref{sec:ft-primitives}
introduces the two reusable redundancy primitives. Section~\ref{sec:ft-wrappers}
shows how they wrap the functional units. Section~\ref{sec:ft-scenarioB}
describes per-functional-unit and full-accelerator hardening (Scenario~A), and
Section~\ref{sec:ft-scenarioA} describes replication-based redundancy
(Scenario~B). Section~\ref{sec:ft-control} explains how a configuration is
selected at elaboration time.

\section{Reusable redundancy primitives}
\label{sec:ft-primitives}

The framework rests on two deliberately tiny, technology-independent VHDL
entities. Their smallness is a feature: because each operates on a single bit,
it can be replicated bit-by-bit across an arbitrarily wide datapath, and because
it is purely structural it imposes no constraints on the module it protects.

\subsection{The DMR comparison cell}

The Dual-Modular-Redundancy cell takes the corresponding output bit from two
copies of a module and produces a single result bit together with an error flag.
If the two inputs agree, the result is forwarded and the flag is deasserted; if
they disagree, the flag is asserted to signal a detected error. The complete
entity is shown in Listing~\ref{lst:dmr}.

\begin{lstlisting}[language=myvhdl, caption={The single-bit DMR comparison cell.}, label={lst:dmr}]
entity DMR is
  port(
    in1, in2 : in  std_logic;
    err_flag : out std_logic;
    result   : out std_logic
  );
end entity DMR;

architecture DMR_arch of DMR is
begin
  process(in1, in2)
  begin
    if in1 = in2 then
      result   <= in1;
      err_flag <= '0';  -- no error
    else
      result   <= in1;  -- arbitrary choice
      err_flag <= '1';  -- error detected
    end if;
  end process;
end architecture DMR_arch;
\end{lstlisting}

When the two copies disagree the cell cannot know which is correct, so it
forwards one of them (here \texttt{in1}) and relies on the asserted
\texttt{err\_flag} to trigger a system-level response. The error flags from all
bits of a functional unit are combined so that any single-bit discrepancy raises
the unit-level error indication.

\subsection{The TMR majority voter}

The Triple-Modular-Redundancy cell takes the corresponding output bit from three
copies of a module and outputs the majority value, realising
Equation~\ref{eq:majority}. The complete entity is shown in
Listing~\ref{lst:tmr}.

\begin{lstlisting}[language=myvhdl, caption={The single-bit TMR majority voter.}, label={lst:tmr}]
entity TMR is
  port(
    in1, in2, in3 : in  std_logic;
    result        : out std_logic
  );
end entity TMR;

architecture TMR_arch of TMR is
begin
  result <= (in1 and in2) or (in2 and in3) or (in1 and in3);
end architecture TMR_arch;
\end{lstlisting}

As long as at most one of the three inputs is wrong at any given bit, the
majority output is correct, so the voter \emph{masks} the fault without any need
for a recovery action. Unlike DMR, TMR therefore does not produce an error flag:
correct operation continues transparently.

\section{Wrapping the functional units}
\label{sec:ft-wrappers}

A redundancy \emph{wrapper} is a VHDL entity that exposes exactly the same port
interface as the functional unit it protects, so that it is a drop-in
replacement. Internally, the wrapper instantiates the functional unit two times
(for DMR) or three times (for TMR), drives all copies with identical inputs, and
recombines their outputs bit-by-bit using arrays of the primitive cells of
Section~\ref{sec:ft-primitives}. Figure~\ref{fig:wrappers} illustrates both
arrangements.

\begin{figure}[htbp]
\centering
\begin{subfigure}[t]{0.48\textwidth}
\centering
\begin{tikzpicture}[
  box/.style={draw, rounded corners, minimum width=16mm, minimum height=7mm, font=\scriptsize, fill=orange!12, align=center},
  cell/.style={draw, minimum width=9mm, minimum height=7mm, font=\scriptsize, fill=blue!8, align=center},
  arr/.style={-{Latex[length=1.6mm]}}]
\node (in) {\scriptsize inputs};
\node[box, right=6mm of in, yshift=6mm] (a) {FU copy 1};
\node[box, below=3mm of a] (b) {FU copy 2};
\node[cell, right=8mm of a, yshift=-7mm] (d) {DMR\\array};
\node[right=6mm of d, yshift=3mm, font=\scriptsize] (r) {result};
\node[right=6mm of d, yshift=-4mm, font=\scriptsize] (e) {err\_flag};
\draw[arr] (in) |- (a.west);
\draw[arr] (in) |- (b.west);
\draw[arr] (a.east) -- (d.west);
\draw[arr] (b.east) -- (d.west);
\draw[arr] (d.east) -- (r);
\draw[arr] (d.east) -- (e);
\end{tikzpicture}
\caption{DMR wrapper: two copies, bitwise comparison, error flag.}
\end{subfigure}
\hfill
\begin{subfigure}[t]{0.48\textwidth}
\centering
\begin{tikzpicture}[
  box/.style={draw, rounded corners, minimum width=16mm, minimum height=6mm, font=\scriptsize, fill=orange!12, align=center},
  cell/.style={draw, minimum width=9mm, minimum height=7mm, font=\scriptsize, fill=blue!8, align=center},
  arr/.style={-{Latex[length=1.6mm]}}]
\node (in) {\scriptsize inputs};
\node[box, right=6mm of in, yshift=8mm] (a) {FU copy 1};
\node[box, below=2mm of a] (b) {FU copy 2};
\node[box, below=2mm of b] (c) {FU copy 3};
\node[cell, right=8mm of b] (v) {TMR\\voter};
\node[right=6mm of v, font=\scriptsize] (r) {result};
\draw[arr] (in) |- (a.west);
\draw[arr] (in) |- (b.west);
\draw[arr] (in) |- (c.west);
\draw[arr] (a.east) -- (v.west);
\draw[arr] (b.east) -- (v.west);
\draw[arr] (c.east) -- (v.west);
\draw[arr] (v.east) -- (r);
\end{tikzpicture}
\caption{TMR wrapper: three copies, bitwise majority vote.}
\end{subfigure}
\caption{Redundancy wrappers. Each preserves the protected unit's interface, so
hardening is transparent to the rest of the accelerator.}
\label{fig:wrappers}
\end{figure}

Listing~\ref{lst:adder_dmr} sketches the structure of the adder DMR wrapper. The
wrapper declares an internal array type that holds the duplicated outputs, binds
the \texttt{DMR} and \texttt{DSP\_Adders} components, instantiates two adder
copies, and applies the DMR cell across every output bit in a generate loop. The
multiplier, shifter, and comparator wrappers follow exactly the same pattern,
differing only in the width and type of the outputs they recombine.

\begin{lstlisting}[language=myvhdl, caption={Structure of the adder DMR wrapper (abridged).}, label={lst:adder_dmr}]
architecture DSP_Adders_DMR_arch of DSP_Adders_DMR is
  type dmr_type is array (1 downto 0)
       of array_2d_dsp_operands(FU_NUM-1 downto 0);
  signal dmr_dsp_out_adder_results : dmr_type;
  component DMR
    port(in1, in2 : in std_logic; err_flag, result : out std_logic);
  end component;
  component DSP_Adders generic(...); port(...); end component;
begin
  -- two identical copies of the adder functional unit
  ADDER_COPY_0 : DSP_Adders generic map(...) port map(...,
      dsp_out_adder_results => dmr_dsp_out_adder_results(0), ...);
  ADDER_COPY_1 : DSP_Adders generic map(...) port map(...,
      dsp_out_adder_results => dmr_dsp_out_adder_results(1), ...);
  -- bitwise comparison across the whole result vector
  GEN_DMR : for i in ... generate
    DMR_i : DMR port map(
      in1 => dmr_dsp_out_adder_results(0)(...)(i),
      in2 => dmr_dsp_out_adder_results(1)(...)(i),
      result => dsp_out_adder_results(...)(i),
      err_flag => err_flag_vector(i));
  end generate;
end architecture;
\end{lstlisting}

\section{Scenario A: per-functional-unit and full-accelerator hardening}
\label{sec:ft-scenarioB}

In the single-lane accelerator, redundancy is applied by replacing one or more
functional units with their hardened wrappers. The top-level design exposes a
set of Boolean constants --- \texttt{ADDERS\_DMR\_EN}, \texttt{ADDERS\_TMR\_EN},
\texttt{MULTIPLIERS\_DMR\_EN}, \texttt{MULTIPLIERS\_TMR\_EN},
\texttt{SHIFTERS\_DMR\_EN}, \texttt{SHIFTERS\_TMR\_EN},
\texttt{COMPARATORS\_DMR\_EN}, and \texttt{COMPARATORS\_TMR\_EN} --- each of which
selects whether the corresponding unit is instantiated as plain, DMR-wrapped, or
TMR-wrapped. This yields the following configurations, all studied
experimentally:

\begin{itemize}
  \item \textbf{Baseline} --- no constant set; the accelerator is unhardened and
        provides the reference area, power, and timing.
  \item \textbf{Single-FU DMR / TMR} --- exactly one of the four arithmetic
        units (adder, multiplier, shifter, or comparator) is wrapped, the others
        left plain. This isolates the cost of protecting each unit individually.
  \item \textbf{Full DMR / TMR} --- all four arithmetic units are wrapped
        simultaneously, giving the cost of protecting the entire arithmetic
        datapath.
\end{itemize}

Because the four functional units differ greatly in size, single-FU hardening
spans a wide range of costs: protecting the comparator is nearly free, whereas
protecting the multiplier is the dominant contributor to any scheme that
includes it. These differences are quantified in Chapter~\ref{ch:results} and
form the basis of the design guidelines.

\section{Scenario B: replication-based redundancy}
\label{sec:ft-scenarioA}

The second way of obtaining redundancy exploits the accelerator's existing
\emph{performance} replication. With \texttt{replicate\_accl\_en} enabled, three
complete accelerator lanes operate in parallel. Rather than building fresh
copies of each functional unit, the redundancy logic is placed in the
\texttt{DSP\_Mapping} layer, where the results produced by the lanes are already
gathered:

\begin{itemize}
  \item \textbf{Replicated baseline} --- three lanes, results simply mapped out,
        no redundancy logic. This is the throughput-oriented configuration.
  \item \textbf{Replicated + DMR voting} --- the mapping layer compares lane
        outputs and raises an error flag on disagreement.
  \item \textbf{Replicated + TMR voting} --- the mapping layer majority-votes the
        three lane outputs, masking any single faulty lane.
\end{itemize}

The mapping unit toggles between these through its own \texttt{DMR\_ENABLED} and
\texttt{TMR\_ENABLED} constants. The crucial observation --- borne out by the
results --- is that since the costly datapath copies already exist for
throughput, the \emph{incremental} cost of adding voting is only the cost of the
comparators or voters in the mapping layer, which is small relative to the
replicated datapath.

\section{Selecting a configuration at elaboration}
\label{sec:ft-control}

A complete experimental configuration is therefore defined by two groups of
settings: the top-level \texttt{replicate\_accl\_en} generic (Scenario~B versus
Scenario~A) and, within each scenario, the redundancy-enable constants. Because
all of these are resolved at \emph{elaboration} time, switching configurations
requires no change to the functional logic --- only a change of constants --- and
each configuration elaborates into a distinct, fully synthesisable netlist. The
synthesis sweep of Chapter~\ref{ch:experiment} automates exactly this: it edits
the relevant constants, elaborates, and synthesises, sweeping the full
configuration matrix across all technology nodes. The complete matrix is listed
in Table~\ref{tab:matrix}.

\begin{table}[htbp]
\centering
\caption{Complete configuration matrix explored in this work. Each configuration is synthesised at all three technology nodes, for a total of 42 synthesis runs.}
\label{tab:matrix}
\small
\begin{tabular}{llp{7.5cm}}
\toprule
Scenario & Configuration & Description \\
\midrule
A (single core) & Baseline & Unhardened reference accelerator \\
 & Adder DMR / TMR & Redundancy on the adder functional unit only \\
 & Multiplier DMR / TMR & Redundancy on the multiplier functional unit only \\
 & Shifter DMR / TMR & Redundancy on the shifter functional unit only \\
 & Comparator DMR / TMR & Redundancy on the comparator functional unit only \\
 & Full DMR / TMR & Redundancy applied to all functional units simultaneously \\
B (replicated) & Replicated baseline & Three accelerator lanes, no voting \\
 & Replicated + DMR & Three lanes with duplication-comparison in the mapping unit \\
 & Replicated + TMR & Three lanes with majority voting in the mapping unit \\
\bottomrule
\end{tabular}
\end{table}

